\documentclass[10pt]{article}
\usepackage[a4paper,margin=0.9in]{geometry}
\usepackage[T1]{fontenc}
\usepackage{lmodern}
\usepackage{amsmath,amssymb}
\usepackage{graphicx}
\usepackage{microtype}
\usepackage[colorlinks=true,urlcolor=blue]{hyperref}
\newcommand{\T}{\mathbb T}
\title{A Kronecker-Set Unitary Is Weakly Supercyclic\\
but Not Weakly Stable}
\author{Literature-implied answer to arXiv:2010.13331}
\date{13 August 2026}

\begin{document}
\maketitle
\paragraph{Status.}
Literature-implied full negative answer to the source's broad question.  The
supporting authors did not identify their construction as an answer to this
later question.

\section*{Question and identification}
Kubrusly--Vieira ask on page 5 whether every weakly supercyclic power-bounded
operator is weakly stable.

\begin{figure}[ht]
\centering
\includegraphics[width=.95\textwidth]{figures/open_problem_crop.png}
\caption{The question in arXiv:2010.13331, page 5.}
\end{figure}

Bayart--Matheron Example 3.6 supplies the counterexample.  Let $E\subset\T$
be a perfect Kronecker set, let $\mu$ be a continuous probability measure
supported on $E$, and let
\[
 U=M_z:L^2(\mu)\to L^2(\mu),\qquad (Uf)(z)=zf(z).
\]
They prove that $\{az^n:a\in\mathbb C,n\ge0\}$ is weakly dense in
$L^2(\mu)$, so $1$ is a weakly supercyclic vector for the unitary $U$.

The same displayed construction says that a Kronecker set is a Dirichlet set:
there are positive integers $p_k\to\infty$ such that
\[
 \sup_{z\in E}|z^{p_k}-1|\longrightarrow0.
\]
Consequently
\[
 \langle U^{p_k}1,1\rangle
   =\int_E z^{p_k}\,d\mu(z)\longrightarrow1,
\]
so $U^n1$ does not converge weakly to $0$.  Thus $U$ is unitary (hence power
bounded), weakly supercyclic, and not weakly stable.

\begin{figure}[ht]
\centering
\includegraphics[width=.95\textwidth]{figures/supporting_answer_crop.png}
\caption{Bayart--Matheron, page 9: weak density and the Dirichlet subsequence
appear in the same construction.}
\end{figure}

\section*{Scope and search evidence}
This implication answers the broad weak-density question negatively.  It does
not make the example weakly $l$-sequentially supercyclic: a weakly convergent
sequence from a fixed projective unitary orbit has bounded scalar multipliers,
which is a materially stronger restriction than weak density.  The companion
full-result packet proves that no unitary of dimension greater than one has
that direct-sequence property.

A bounded search through 13 August 2026 checked the run indexes, exact wording,
the 2024 survey arXiv:2405.02752, arXiv:2306.08197, and the cited
Bayart--Matheron and Shkarin papers.  Later sources continue to call the broad
question open; no source found explicitly records the two-line implication
above.  This packet therefore records an agent-identified literature
implication, not a claim of a new construction.

\paragraph{Human-review recommendation.}
Check only the sign convention in the displayed inner product and confirm that
the positive Dirichlet sequence $p_k$ in Bayart--Matheron is the same sequence
shown in the supporting excerpt.  Either convention gives convergence to
$1$ after complex conjugation.

\begin{thebibliography}{9}
\bibitem{KV}
C. S. Kubrusly and P. C. M. Vieira,
\emph{Boundedly Spaced Subsequences and Weak Dynamics},
Journal of Function Spaces (2018), Article ID 4732836;
\href{https://arxiv.org/abs/2010.13331}{arXiv:2010.13331}.
\bibitem{BM}
F. Bayart and E. Matheron,
\emph{Hyponormal operators, weighted shifts and weak forms of
supercyclicity}, Proc. Edinburgh Math. Soc. 49 (2006), 1--15,
doi:10.1017/S0013091504000975.
\end{thebibliography}
\end{document}
